\documentclass[11pt,a4paper]{article}
\usepackage[utf8]{inputenc}
\usepackage[T1]{fontenc}
\usepackage[margin=0.72in]{geometry}
\usepackage{enumitem}
\usepackage{titlesec}
\usepackage{xcolor}
\usepackage{hyperref}
\usepackage{microtype}

\definecolor{primarycolor}{RGB}{33,37,41}
\definecolor{secondarycolor}{RGB}{0,92,150}
\definecolor{accentcolor}{RGB}{90,100,110}
\hypersetup{colorlinks=true,urlcolor=secondarycolor,linkcolor=secondarycolor}
\titleformat{\section}{\large\bfseries\color{primarycolor}}{}{0em}{}
\titlespacing*{\section}{0pt}{1.0em}{0.35em}
\setlist[itemize]{leftmargin=1.3em,itemsep=0.12em,topsep=0.2em,parsep=0pt}
\pagestyle{empty}
\setlength{\parindent}{0pt}
\setlength{\parskip}{0.25em}

\begin{document}
\begin{center}
{\LARGE\bfseries Marta Kowal-Brandt}\\[0.25em]
{\large DevOps / Platform Engineering | CI/CD, Kubernetes and Infrastructure Automation}\\[0.25em]
Wroclaw, Poland | \href{mailto:marta.kowalbrandt@example.com}{marta.kowalbrandt@example.com} | +48 555 303 303 | \href{https://gitlab.com/mkb-example}{gitlab.com/mkb-example}
\end{center}

\section*{PROFESSIONAL SUMMARY}
QA automation engineer with 7 years of experience and a strong focus on CI/CD, test infrastructure and build tooling. Owns GitLab CI pipelines and Kubernetes-based ephemeral test environments for three product teams. Experienced in improving delivery pipelines, automating environment lifecycle tasks, and making test systems more observable. Building further infrastructure experience through a personal k3s and GitOps homelab and Terraform practice.

\section*{SELECTED INFRASTRUCTURE PROJECTS}
\textbf{Merge-Request Ephemeral Environments | Odrava Health Software | 2023–2024}
\begin{itemize}
\item Built short-lived environments for each merge request, deploying the application, PostgreSQL and seeded data to a shared Kubernetes cluster using Helm from GitLab CI.
\item Automated environment cleanup after merge or 48 hours of inactivity; enabled QA and product reviews before merge and e2e tests against each environment.
\item Technologies: GitLab CI, Helm, Kubernetes, Docker, Bash, PostgreSQL.
\end{itemize}

\textbf{CI Pipeline Speed-up | Odrava Health Software | 2022}
\begin{itemize}
\item Profiled pipelines across three teams, parallelised test jobs, added dependency caching and separated e2e testing into its own stage.
\item Reduced average pipeline time from 38 to 14 minutes and monthly CI runner cost by about 30\%.
\item Technologies: GitLab CI, Docker, pytest-xdist.
\end{itemize}

\textbf{Homelab k3s Cluster with GitOps | Personal | 2024–2025}
\begin{itemize}
\item Operated a three-node Raspberry Pi 5 k3s cluster, deploying applications from a Git repository with Argo CD.
\item Set up Prometheus, Grafana and Loki monitoring, Traefik ingress and cert-manager for TLS.
\item Provisioned infrastructure for a small Hetzner VPS using Terraform; used Ansible at a basic level.
\item Technologies: k3s, Argo CD, Prometheus, Grafana, Loki, Terraform, Ansible.
\end{itemize}

\textbf{Flaky Test Dashboard | Odrava Health Software | 2023}
\begin{itemize}
\item Collected CI test results in PostgreSQL and created Grafana dashboards showing flaky tests by team.
\item The flaky test rate fell from 7\% to below 1\% after teams began addressing the top offenders.
\item Technologies: Python, PostgreSQL, Grafana, GitLab API.
\end{itemize}

\section*{PROFESSIONAL EXPERIENCE}
\textbf{DevOps Engineer | Odrava Health Software, Wroclaw | 2021–Present}
\begin{itemize}
\item Own GitLab CI pipelines for three product teams (approximately 25 developers), covering build, unit, integration and e2e test stages and artifact publishing.
\item Built and maintain merge-request test environments using Docker Compose and Helm on the company’s shared Kubernetes cluster; have namespace-level access.
\item Maintain a pytest and Playwright e2e suite of approximately 900 tests; reduced the flaky test rate from 7\% to below 1\%.
\item Develop Bash and Python tooling for PostgreSQL test-data seeding.
\item Wrote the internal “How our CI works” onboarding document and mentored two QA interns in 2023.
\end{itemize}

\textbf{QA Engineer | Brightmill Software House, Wroclaw | 2019–2021}
\begin{itemize}
\item Tested web and mobile applications for client projects using manual and automated approaches.
\item Introduced a Selenium and Java automated regression suite for the largest client project.
\item Configured Jenkins jobs for nightly regression runs.
\end{itemize}

\textbf{Customer Support Specialist (2nd line) | Brightmill Software House | 2018–2019}
\begin{itemize}
\item Investigated customer issues using logs and databases, and reproduced bugs for developers.
\end{itemize}

\section*{TECHNICAL SKILLS}
CI/CD and containers: GitLab CI, Jenkins, Docker, Docker Compose\\
Infrastructure and orchestration: Kubernetes (namespace-level access to a shared cluster), Helm, k3s, Terraform (personal project), Argo CD, Ansible (basic)\\
Monitoring and data: Prometheus, Grafana, Loki, PostgreSQL, SQL\\
Programming and automation: Python, Bash, Java (test code), Git\\
Cloud exposure: AWS EC2 and S3 fundamentals through personal learning; Hetzner VPS (personal project)\\
Testing: pytest, Playwright, Selenium, Postman\\
Operating systems: Linux

\section*{EDUCATION}
BEng, Computer Engineering | Silesian Lowlands University of Technology | 2018

\section*{CERTIFICATIONS AND PROFESSIONAL DEVELOPMENT}
ISTQB Foundation Level | 2019\\
Certified Kubernetes Administrator (CKA) exam scheduled for January 2027; currently studying

\section*{LANGUAGES}
Polish (native), English (C1), German (A2)
\end{document}